\documentclass[conference]{IEEEtran}

\usepackage{cite}
\usepackage{amsmath,amssymb,amsfonts}
\usepackage{graphicx}
\usepackage{booktabs}
\usepackage{array}
\usepackage{url}
\usepackage{xcolor}

\title{Notes on Multiple Linear Regression}

\author{
\IEEEauthorblockN{Luis Mendez}
}

\begin{document}
\maketitle

\section{Building a model}
\begin{itemize}
    \item Define the problem and collect data.
    \item Preprocess the data (handle missing values, encode categorical variables, etc.).
    \item Split the data into training and testing sets.
    \item Train the linear regression model on the training data.
    \item Evaluate the model using metrics such as Mean Absolute Error (MAE), Mean Squared Error (MSE), or R-squared on the testing data.
    \item Interpret the coefficients to understand the relationship between independent and dependent variables.
\end{itemize}

\section*{Backward Elimination}
\begin{itemize}
    \item Step 1: Select a significance level (e.g., 0.05) to stay in the model.
    \item Step 2: Fit the model and calculate the p-values for each coefficient.
    \item Step 3: Remove the variable with the highest p-value (greater than a significance level, e.g., 0.05).
    \item Step 4: Repeat the process until all remaining variables have p-values below the significance level.
    \item Step 5: The final model will contain only the statistically significant variables. Fit the model without the variables that were removed and evaluate its performance.
\end{itemize}

\section*{Forward Selection}
\begin{itemize}
    \item Step 1: Select a significance level (e.g., 0.05) to enter the model.
    \item Step 2: Fit the model with each independent variable separately and calculate the p-values.
    \item Step 3: Add the variable with the lowest p-value (below the significance level) to the model.
    \item Step 4: Repeat the process by adding variables one at a time, checking the p-values at each step.
    \item Step 5: The final model will contain only the statistically significant variables. Fit the model with the selected variables and evaluate its performance.
\end{itemize}

\section*{Bidirectional Elimination}
\begin{itemize}
    \item Step 1: Select a significance level (e.g., 0.05) for both entering and staying in the model.
    \item Step 2: Start with an empty model and perform forward selection to add variables based on their p-values.
    \item Step 3: After adding a variable, check the p-values of all variables in the model. If any variable has a p-value greater than the significance level, remove it from the model.
    \item Step 4: Repeat the process of adding and removing variables until no more variables can be added or removed based on the significance level.
    \item Step 5: The final model will contain only the statistically significant variables. Fit the model with the selected variables and evaluate its performance.    
\end{itemize}

\section*{All possible models}
\begin{itemize}
    \item Step 1: Generate all possible combinations of independent variables.
    \item Step 2: Fit a linear regression model for each combination of variables and calculate the performance metrics (e.g., R-squared, Adjusted R-squared, AIC, BIC).
    \item Step 3: Compare the performance metrics across all models to identify the best-performing model.
    \item Step 4: Select the model with the best performance metric (e.g., highest Adjusted R-squared or lowest AIC/BIC) as the final model.
    \item Step 5: Fit the final model with the selected variables and evaluate its performance on the testing data.
\end{itemize}

\end{document}